% !TEX root = cv-ko.tex
% !TEX program = xelatex
% Korean CV content. Edit this file for wording, dates, and entries.
\def\cvpdftitle{hy30nq | 보안 연구 CV}
\def\cvfooterlabel{hy30nq · 보안 연구 CV}
\def\cvupdated{2026년 9월 업데이트}
% Shared monochrome resume layout for cv-ko.tex and cv-en.tex.
% Edit typography, spacing, and entry macros in this file.
\documentclass[10pt,a4paper]{article}

\usepackage[a4paper,left=18mm,right=18mm,top=15mm,bottom=15mm]{geometry}
\usepackage{fontspec}
\usepackage{array}
\usepackage{tabularx}
\usepackage{enumitem}
\usepackage{fancyhdr}
\usepackage{hyperref}
\usepackage{xcolor}
\usepackage{ragged2e}
\usepackage{needspace}

% Local macOS preference, followed by portable fallbacks for XeLaTeX.
\IfFontExistsTF{NanumSquare Neo}{
  \setmainfont{NanumSquare Neo}[BoldFont={NanumSquare Neo ExtraBold}]
  \setsansfont{NanumSquare Neo}[BoldFont={NanumSquare Neo ExtraBold}]
}{
  \IfFontExistsTF{Pretendard}{
    \setmainfont{Pretendard}
    \setsansfont{Pretendard}
  }{
    \IfFontExistsTF{Noto Sans CJK KR}{
      \setmainfont{Noto Sans CJK KR}
      \setsansfont{Noto Sans CJK KR}
    }{
      \setmainfont{Arial Unicode MS}
      \setsansfont{Arial Unicode MS}
    }
  }
}

% Strict black-and-white palette.
\definecolor{ink}{HTML}{000000}

\hypersetup{
  unicode=true,
  hidelinks,
  pdftitle={\cvpdftitle},
  pdfauthor={hy30nq},
  pdfsubject={Security research, vulnerability analysis, and CTF engineering}
}

\pagestyle{fancy}
\fancyhf{}
\renewcommand{\headrulewidth}{0pt}
\renewcommand{\footrulewidth}{0pt}
\fancyfoot[C]{%
  \makebox[\textwidth]{%
    \fontsize{7.2}{9}\selectfont
    \cvfooterlabel\hfill\cvupdated\ · \thepage
  }%
}
\setlength{\footskip}{8mm}

\setlength{\parindent}{0pt}
\setlength{\parskip}{0pt}
\setlength{\tabcolsep}{0pt}
\renewcommand{\arraystretch}{1.0}
\hyphenpenalty=10000
\exhyphenpenalty=10000
\color{ink}

\setlist[itemize]{
  leftmargin=4.2mm,
  itemsep=1.6pt,
  topsep=2.3pt,
  parsep=0pt,
  partopsep=0pt,
  label=\raisebox{0.15ex}{\tiny\textbullet}
}

\newcommand{\cvheader}[1]{%
  \begin{tabularx}{\textwidth}{@{}X>{\RaggedLeft\arraybackslash}p{66mm}@{}}
    {\fontsize{30}{32}\selectfont\bfseries hy30nq} &
    \begin{tabular}[t]{@{}r@{}}
      \fontsize{8.15}{10.6}\selectfont hy30nq@gmail.com\\
      \fontsize{8.15}{10.6}\selectfont github.com/hy30nq03\\
      \fontsize{8.15}{10.6}\selectfont hyeonql.tistory.com
    \end{tabular}\\[-0.2mm]
    {\fontsize{11.2}{13.6}\selectfont\bfseries #1} &
    {\fontsize{8.15}{10.6}\selectfont hy30nq.com}\\
  \end{tabularx}
  \vspace{2.6mm}
  \hrule height 0.8pt
}

\newcommand{\sectiontitle}[2]{%
  \needspace{18mm}
  \vspace{5.2mm}%
  \begin{tabularx}{\textwidth}{@{}Xr@{}}
    {\fontsize{12.8}{15}\selectfont\bfseries #1} &
    {\fontsize{7.2}{9}\selectfont\bfseries\MakeUppercase{#2}}
  \end{tabularx}
  \vspace{1.1mm}
  \hrule height 0.45pt
  \vspace{2.8mm}
}

\newcommand{\entry}[4]{%
  \needspace{14mm}
  \begin{tabularx}{\textwidth}{@{}>{\RaggedRight\arraybackslash}p{24mm}>{\RaggedRight\arraybackslash}X@{}}
    {\fontsize{8.1}{10.2}\selectfont #1} &
    {\fontsize{9.8}{12.2}\selectfont\bfseries #2}\\[-0.2mm]
    & {\fontsize{8.0}{10.2}\selectfont\itshape #3}\\[0.55mm]
    & {\fontsize{8.65}{11.3}\selectfont #4}
  \end{tabularx}
  \vspace{3.1mm}
}

\newcommand{\compactentry}[3]{%
  \needspace{6mm}
  \begin{tabularx}{\textwidth}{@{}>{\RaggedRight\arraybackslash}p{24mm}>{\RaggedRight\arraybackslash}X@{}}
    {\fontsize{8.1}{10.2}\selectfont #1} &
    {\fontsize{8.75}{11.1}\selectfont #2\enspace
     {\fontsize{7.9}{9.8}\selectfont · #3}}
  \end{tabularx}
  \vspace{1.7mm}
}

% Compact bullet list for CTF authoring and participation history.
\newenvironment{cvlist}{%
  \begin{itemize}[
    leftmargin=4.2mm,
    itemsep=2.0pt,
    topsep=0.5pt,
    parsep=0pt,
    partopsep=0pt
  ]
  \fontsize{8.75}{11.1}\selectfont
}{%
  \end{itemize}
}

\newcommand{\tag}[1]{%
  \begingroup
  \setlength{\fboxsep}{1.2mm}%
  \setlength{\fboxrule}{0.4pt}%
  \fbox{\strut\fontsize{7.7}{9.5}\selectfont\bfseries #1}%
  \endgroup
}

\newcommand{\cvskills}{%
  \vfill
  \begin{center}
    \tag{Vulnerability Research}\hspace{1.4mm}
    \tag{CTF Engineering}\hspace{1.4mm}
    \tag{Security Education}\hspace{1.4mm}
    \tag{Python · Docker · Git}
  \end{center}
}


\begin{document}
\sffamily

\cvheader{Security Researcher · CTF Builder}

% =============================================================================
% 소속
% =============================================================================
\sectiontitle{소속}{Affiliations}

\compactentry{2023--현재}
  {\textbf{고려대학교 인공지능사이버보안학과}}
  {개인정보보호융합전공 · 학·석사 연계과정생}

\compactentry{현재}
  {\textbf{과학기술전문사관 제3기 후보생}}
  {Science \& Technology Specialist Officer Program}

\compactentry{}
  {\textbf{SeeCure Lab · KUality · Knights of the SPACE · Rubiyalab}}
  {Security Research \& CTF Communities}

% =============================================================================
% 수상: 이 섹션에만 링크 사용
% =============================================================================
\sectiontitle{수상}{Awards}

\compactentry{2025}
  {\href{https://security.korea.ac.kr/m06_04/59}{\textbf{장려상} · 제2회 개인정보보호·보안 정책 아이디어 공모전}}
  {고려대학교 정보보호대학원}

\compactentry{2024}
  {\href{https://blog.naver.com/kisiaedu/223785868393}{\textbf{우수상 · 한국정보보호산업협회장상} · AI보안 기술개발 참여수기 공모전}}
  {한국정보보호산업협회}

\compactentry{2024}
  {\href{https://www.kisia.or.kr/talent_support/education_business/ai_security_technology_development/}{\textbf{과학기술정보통신부 장관상 · 최우수 수료생} · KISIA AI보안 기술개발 교육과정}}
  {교육·실습·프로젝트 종합 평가}

\compactentry{2024}
  {\href{https://www.boannews.com/media/view.asp?idx=132213\&direct=mobile}{\textbf{1등 · 과학기술정보통신부 장관상} · KISIA 정보보호 개발자 해커톤}}
  {크롬 확장 프로그램 기반 악성 URL 탐지 서비스}

\compactentry{2024}
  {\href{https://www.msit.go.kr/bbs/view.do?sCode=user\&mId=113\&mPid=238\&bbsSeqNo=94\&nttSeqNo=3184227}{\textbf{Top 20 · 과학기술정보통신부 장관 인증서} · 화이트햇 스쿨 1기}}
  {320명 중 상위 20인}

% =============================================================================
% CTF 본선 참가
% =============================================================================
\sectiontitle{CTF 본선 참가}{CTF Finals}

\begin{cvlist}
  \item \textbf{2026 · DEF CON 34 CTF} · 본선 진출 · Team SSS
  \item \textbf{2025 · 화이트햇콘테스트} · 본선 진출 · 예선 5위 · Team KUality
  \item \textbf{2025 · Black Hat MEA CTF} · 본선 진출 · 예선 41위 · Team PPU-AENG Overflow
\end{cvlist}

% =============================================================================
% CTF 출제: 항목은 최신순으로 한 줄씩 추가
% =============================================================================
\sectiontitle{CTF 출제}{Challenge Authoring \& Operations}

\begin{cvlist}
  \item \textbf{2026 · KCTF 2026 (KUality)} · 문제 출제 · Liberation Signal 1--3
  \item \textbf{2026 · CSKWS 연합 세미나 CTF} · 대회 운영 및 문제 출제 · Ghost Display, Black Cache
  \item \textbf{2026 · 핵테온 세종 초·중·고 사이버보안 한마당} · CTF 운영 및 문제 출제
  \item \textbf{2026.08 · 정보보호 실무기초 교육과정} · 기초 과정 교육 및 Mini CTF 운영
  \item \textbf{2026.02 · INC0GNITO FESTIVAL CTF} · 대회 운영 및 예선·본선 문제 출제 · Project Incoherence
  \item \textbf{2026.02 · 제5회 한양 해킹방어대회 HCTF with HSPACE} · 문제 출제
  \item \textbf{2025.08 · SSD CTF} · 문제 출제
\end{cvlist}

\newpage

% =============================================================================
% 프로젝트
% =============================================================================
\sectiontitle{프로젝트}{Selected Projects}

\compactentry{2026--현재}
  {\textbf{AI 시스템 대상 적대적 공격·사고 추적 기술 연구}}
  {참여 연구원 · 수행 중}

\compactentry{2026--현재}
  {\textbf{원격 장비관리 통신 프로토콜 보안성 분석(TR-069) II}}
  {참여 연구원 · 수행 중}

\compactentry{2026--현재}
  {\textbf{Threat Actor 캠페인 기반 미래 위협 시나리오 예측 AI 기술 및 MITRE D3FEND 기반 대응 기법 고도화}}
  {참여 연구원 · 수행 중}

\entry{2026}
  {AI IDE 보안 강화를 위한 Tool/MCP 호출 게이트웨이 기반 방어 방법론}
  {프로젝트 학기제 · Vulnerability Research}
  {Cursor IDE의 비신뢰 프로젝트 Hook 설정이 워크스페이스 신뢰 검증과 샌드박스 없이 셸 명령을 실행하는 문제를 분석하고 재현했습니다.
   재현 가능한 PoC와 책임 있는 제보 문서를 정리하고, Verify·Restrict·Log 게이트웨이 관점의 방어 방법론을 설계했습니다.}

\entry{2025.06--07}
  {HSPACE 침투 테스트 교육 콘텐츠}
  {Content Developer · Knights of the SPACE}
  {입문자가 단계적으로 공격 흐름을 이해할 수 있도록 가상 쇼핑몰과 난이도별 침투 테스트 시나리오를 구성하고,
   반복 실행 가능한 실습 환경을 제작했습니다.}

\entry{2025.06--08}
  {실전 버그바운티 사례 기반 CTF 플랫폼}
  {Backend Developer · FinderGap}
  {실제 버그바운티 사례를 재현 가능한 CTF 문제로 전환하는 교육 플랫폼의 백엔드를 구현했습니다.}

\entry{2024}
  {WordPress 플러그인 취약점 분석 자동화}
  {고려대학교 인공지능사이버보안학과 학술제}
  {WordPress 플러그인을 대상으로 한 자동 취약점 탐지 도구를 공동 개발하고 SQL Injection 분석 모듈을 담당했습니다.
   분석 결과와 구현 내용을 교내 학술제에서 발표했습니다.}

\entry{2024}
  {AI 기반 악성 URL 분류기}
  {Frontend · Data Preprocessing · KISIA AI Security Network}
  {피싱 URL 데이터 정제와 모델 학습 파이프라인 구축에 참여하고,
   크롬 확장 프로그램 형태의 악성 URL 탐지 서비스 프론트엔드를 개발했습니다.}

% =============================================================================
% CVE
% =============================================================================
\sectiontitle{CVE}{Vulnerability Disclosure}

\entry{2024}
  {CVE-2024-11950 · XnSoft XnView Classic RWZ 정수 언더플로우}
  {Independent Vulnerability Research · ZDI-24-1640 · High Severity}
  {조작된 RWZ 파일 처리 과정의 정수 언더플로우와 메모리 손상을 분석·재현하고 책임 있게 보고했습니다.
   해당 취약점은 원격 코드 실행 영향으로 ZDI 권고문과 CVE로 공개되었습니다.}

% =============================================================================
% 기타
% =============================================================================
\sectiontitle{기타}{Other}

\entry{2025.03--08}
  {KITRI 화이트햇 스쿨 3기 조교}
  {Stage 1 · Midterm · Stage 3}
  {교육 단계별 실습과 중간발표 운영을 지원하고 학습자 피드백과 진행 관리를 담당했습니다.}

\entry{2024.06--현재}
  {KnockOn Q\&A 멘토}
  {Security Community Mentor}
  {웹 보안과 취약점 분석 학습자의 질문에 답변하고 실습 방향을 안내하고 있습니다.}

\compactentry{2024}
  {\textbf{정보처리기능사}}
  {국가기술자격}

\compactentry{}
  {\textbf{언어} · 한국어(모국어) · 영어(읽기·쓰기) · 중국어(HSK 4급, HSKK 초급)}
  {Languages}

\cvskills

\end{document}
